\scanpage{67}
\mspage{63}
% source lines 1--6
\[
\bar\alpha(\beta(x'))\mid\circ\mid\bar\alpha(\beta(y'))
\]
(\unsure{comp.} de $\alpha$ aux disjonctions), pour en conclure, comme on voudrait, $x'\mid\circ\mid y'$, il faudrait que l'inclusion dans $(*)$ soit une identité; ou plutôt que
\[
\overset{\deleted{\unsure{E}}}{\operatorname{supp}}^{\circ}_{E'}(\bar\alpha(\beta(x'))=\operatorname{supp}^{\circ}_{E'}(\{x'\}) .
\]
% source lines 7--11
Voici un choix \uline{optimal} de $\beta$, \unsure{en termes} de $\alpha$:

Soit $\beta=\deleted{\unsure{t}}\alpha^*$ défini par
\[
\deleted{\unsure{\beta}}\alpha^*(x')=\{x\in E\mid\alpha(x)\subset\operatorname{supp}^{\circ}_{E'}(\{x'\})\}
\]
% source lines 12--16: abandoned formulations
\deleted{Notre hypothèse supplémentaire sur}
\deleted{Avec cette \unsure{définition}\inserted{ce choix de $\beta$}, on aura}
\deleted{Avec cette}
\unsure{On aura} alors tautologiquement
\[
\bar\alpha(\alpha^*(x'))\subset\operatorname{supp}^{\circ}_{E'}(\{x'\}),
\]
% source lines 17--20
\deleted{et il faut supposer que cette \unsure{inclusion|identité}}
\deleted{\unsure{de fait|du fait}}\inserted{\unsure{donc}}
\[
\operatorname{supp}^{\circ}_{E'}(\bar\alpha(\alpha^*(x')))\subset\operatorname{supp}^{\circ}_{E'}(\{x'\})
\]
et il faut supposer que cette inclusion est une \uline{égalité}
\[
\boxed{\operatorname{supp}^{\circ}_{E'}(\bar\alpha(\alpha^{\unsure{*|t}}(x')))=\operatorname{supp}^{\circ}_{E'}(\{x'\}).}\ .\tag*{(i)}
\]
% source lines 21--25
(\unsure{Comparer} l'hypothèse de surjectivité dans l'exemple de \unsure{tantôt}). De plus, la propriété que $\alpha$ soit compatible, doit être remplacée par une hyp. de ``compatibilité fidèle''
\[
\boxed{\text{Si }x,y\in E,\quad x\mid\circ\mid y\Longleftrightarrow\bar\alpha(x)\mid\circ\mid\bar\alpha(y)}\tag*{(ii)}
\]
% archive sheet number: 66
